\section{Por qué la exploración es más difícil en Meta-RL}

\subsection{El problema de exploración en Meta-RL no es igual a la exploración ordinaria}

En el RL ordinario, el objetivo de la exploración es encontrar acciones con alto retorno; en Meta-RL, en cambio, la exploración asume además otra responsabilidad: \textbf{identificar cuál es realmente la tarea actual}. Por ello, una misma conducta puede tener valor a largo plazo aunque no produzca reward en el corto plazo, si reduce de forma considerable la incertidumbre sobre la tarea.

\begin{figure}[H]
\centering
\includegraphics[width=0.9\textwidth]{slides-images/slide-29.jpg}
\caption{La solución más directa: aprender de extremo a extremo exploration y exploitation al mismo tiempo}
\end{figure}

\subsection{Ejemplo del pasillo: exploración que aporta información pero no recompensa}

La clase explicó esto con el ejemplo hallway. Si para saber en qué extremo está la meta primero hay que ver la indicación en la pared, entonces «ir primero a ver la indicación» es la conducta de meta-exploration correcta, pero por sí misma no produce directamente recompensa de la tarea. Así, cuando el aprendizaje de extremo a extremo optimiza solo según el reward final, al optimizador le cuesta reconocer el valor a largo plazo de esta exploración inicial sin recompensa.

\begin{figure}[H]
\centering
\includegraphics[width=0.9\textwidth]{slides-images/slide-31.jpg}
\caption{Ejemplo del pasillo: la conducta de exploración realmente útil no necesariamente produce reward de inmediato}
\end{figure}

\begin{warningbox}{Problema de acoplamiento (coupling problem)}
Si la exploración no se hace bien, la fase de ejecución no obtiene un buen reward; y si la fase de ejecución no obtiene un buen reward durante mucho tiempo, la fase de exploración tampoco recibe una señal de entrenamiento suficientemente clara. Esta dependencia mutua puede hacer que la optimización quede atrapada en un óptimo local muy malo.
\end{warningbox}

\subsection{Ejemplo de la cocina: exploración y ejecución se bloquean mutuamente}

El ejemplo de cooking que citó el ponente deja más claro este problema: si el agent no encuentra primero los ingredientes, no puede aprender a cocinar; pero si nunca aprende a cocinar, el reward final no basta para indicarle qué forma de buscar ingredientes es valiosa. El resultado es que exploration y execution se frenan mutuamente.

\begin{figure}[H]
\centering
\includegraphics[width=0.9\textwidth]{slides-images/slide-33.jpg}
\caption{El coupling problem en el ejemplo de la cocina: la exploración y la ejecución dependen una de otra, lo que dificulta el entrenamiento}
\end{figure}

\subsection{Enfoque alternativo: incorporar explícitamente la estructura de la exploración}

Por ello, la clase no se quedó en «basta con aprender de extremo a extremo», sino que presentó además varios tipos de estrategias alternativas:
\begin{enumerate}
    \item \textbf{Posterior sampling / Thompson sampling}: mantener explícitamente la posterior sobre la tarea, muestrear de ella una tarea latente y actuar en consecuencia.
    \item \textbf{Intrinsic rewards}: otorgar una recompensa adicional por la recolección de información en sí misma.
    \item \textbf{Task dynamics / reward prediction}: aprender explícitamente un modelo de la dinámica o de la recompensa relacionado con la tarea, para ayudar a distinguir las tareas.
\end{enumerate}

\begin{figure}[H]
\centering
\includegraphics[width=0.9\textwidth]{slides-images/slide-35.jpg}
\caption{Solution 2a: PEARL realiza posterior sampling mediante una latent task posterior}
\end{figure}

\begin{figure}[H]
\centering
\includegraphics[width=0.9\textwidth]{slides-images/slide-36.jpg}
\caption{Ventajas y desventajas de las estrategias de exploración alternativas: son más fáciles de optimizar, pero en algunos entornos pueden ser claramente subóptimas}
\end{figure}

\begin{importantbox}{La intuición central de PEARL}
PEARL aprende la distribución a priori y la posterior de una latent task variable $z$, y luego condiciona la política a $z$. Una vez que la posterior se concentra, la política ya no necesita «adivinar» la tarea actual a partir del estado oculto, sino que actúa explícitamente con base en la belief actual. Esto tiene más estructura que una red de memoria puramente de caja negra, pero también introduce supuestos de modelado más fuertes.
\end{importantbox}
